\section{Síntesi i posicionament del TFM}

La revisió mostra que els components conceptuals del projecte estan establerts,
però apareixen dispersos. La teoria de cues aporta mètriques i criteris
d'equitat; els estudis de planificació descriuen els compromisos entre prioritat
i espera; els marcs de simulació faciliten repetir experiments; i els brokers
mostren les restriccions d'una implementació real. La contribució prevista no
substitueix aquestes línies, sinó que n'integra una part concreta en un prototip
web coherent.

El motor representarà tasques asíncrones no preemptives, workers limitats i un
rellotge virtual. Les polítiques s'implementaran darrere d'una interfície comuna
per executar FIFO, prioritat estricta i prioritat amb envelliment sobre la
mateixa entrada. La persistència conservarà els resultats i les dades necessàries
per reproduir-los. La interfície permetrà preparar escenaris i comparar
mètriques sense inspeccionar logs o taules internes.

Aquest posicionament és més reduït que GridSim, CloudSim o SimGrid de manera
intencionada. No es modelaran topologies de xarxa, màquines virtuals ni
negociació de recursos, ni es reproduiran totes les garanties d'un broker.
L'objectiu és fer visible la relació entre càrrega, política i espera amb un
model comprensible i verificable.

La plataforma s'adreça principalment a estudiants i professionals de backend
que vulguin explorar una decisió abans d'aplicar-la a un sistema real. Una
simulació sintètica pot revelar comportaments, detectar riscos i comparar
alternatives dins del model, però no pot garantir l'efecte d'una política sobre
una organització sense calibrar-la amb dades representatives.
